\documentclass[11pt]{article}
\usepackage[T1]{fontenc}
\usepackage[utf8]{inputenc}
\usepackage[margin=1in]{geometry}
\usepackage{amsmath,amssymb}
\usepackage{hyperref}
\hypersetup{hidelinks}
\setlength{\emergencystretch}{3em}
\setlength{\parindent}{0pt}
\setlength{\parskip}{6pt}
\title{P02.2: fixed-neighborhood LFS regularization in a proper-fibre sector, and a boundary obstruction}
\author{Unverified research draft}
\date{4 October 2026}
\begin{document}
\maketitle

This is an unverified partial mathematical argument. The original characterization of all higher-dimensional good strict analytic spaces remains open here. Only P02.2 is addressed. No unpublished book assertion is used. The results below concern arbitrary classical psh functions, including discontinuous functions and the value $-\infty$.

\section*{1. Precise results and scope}

Let $K$ be complete, algebraically closed, and nontrivially non-Archimedean valued. Let $C$ be an ordinary open subset of the analytification of a smooth connected algebraic $K$-curve. Its genus is unrestricted. The curve is boundaryless. Write $\Gamma=-\log|K^\times|$; it is divisible and dense in $\mathbb R$.

Curve result. For every classical psh function $g:C\to\mathbb R\cup\{-\infty\}$ and every $c\in C$, there is one ordinary open neighborhood $U$ of $c$ and a decreasing net of finite strict-domain real-PL LFS functions on $U$ with pointwise infimum $g|_U$. If $K$ has a countable dense subset in its valuation topology, this net has a decreasing sequence extraction with the same pointwise limit. No continuity of $g$ is assumed. No restriction on the residue field or the residue genera is imposed. The approximants are actual LFS functions in GJR Definition 7.8, rather than a further limiting stage.

Higher-dimensional result. Under these field hypotheses, let $f:X\to C$ be either the actual analytic projective bundle $\mathbb P(E)$ of a locally free analytic sheaf $E$ of rank $r+1$, $r\geq1$, or the product projection $C\times P^{\mathrm{an}}\to C$, where $P$ is a nonempty geometrically connected proper algebraic $K$-variety of positive dimension. Every arbitrary classical psh function on $X$, or on a saturated ordinary open set $f^{-1}(C_0)$, has the preceding decreasing LFS net near each point, on a fixed saturated ordinary neighborhood $f^{-1}(U)$. It has a sequence when $K$ is countably dense. These are higher-dimensional sectors with no torus or invariance hypothesis on the starting function. Properness forces fibre constancy; it is not an imposed restriction on the function.

Boundary result. If $K$ is algebraically closed with uncountable residue field, the closed unit disc has a classical subharmonic function which is not a pointwise limit of a decreasing sequence of continuous functions on any ordinary neighborhood of its Gauss point. Thus it has no decreasing LFS sequence there. The same obstruction occurs on the good strict spaces $D\times\mathbb P^r_K{}^{\mathrm{an}}$, $r\geq1$. For this particular disc function, countability of the residue field is the sharp condition: an explicit decreasing LFS sequence exists when the residue field is countable. This is a sequence obstruction, not a proof of failure for nets.

The higher-dimensional positive result quantifies over all classical psh functions defined on the displayed saturated spaces. It does not establish regularization for arbitrary germs defined only on a nonsaturated small open subset of a projective bundle. A local germ need not extend to a saturated neighborhood. In particular this sector must not be read as a characterization of the full psh sheaf on every open subset of $X$.

\section*{2. Published foundations and exact conventions}

GJR denotes the cached paper arXiv:2111.05741; Thu denotes Thuillier's thesis tel-00010990v1; GRc denotes arXiv:2506.13548v2; GRpl denotes the cached arXiv:2503.05513v1. GRc Definition 4.1 tests pullbacks on analytic curves after every analytic field extension, with upper semicontinuity and values in $\mathbb R\cup\{-\infty\}$. On a smooth boundaryless strict curve this is Thuillier subharmonicity, allowing an identically $-\infty$ connected component: GRc 3.4.1, 4.1.2 and Proposition 3.5. Scalar extension, pullback and decreasing nets preserve this class, by GRc Propositions 3.5 and 4.2.

GJR Definition 5.1 uses the G-topology generated by strictly affinoid domains: a real-PL function is G-locally a finite real linear combination of logarithms of analytic units. Thu 2.1.1 distinguishes the all-domain site $X_G$ from the strict rigid site $X_R$, and Thu Definition 3.2.1 uses all affinoid domains to define its lisse functions. These are not silently identified. GJR Example 5.3 explicitly shows that an irrational bend of a maximum of real-PL functions can fail to be strict-domain real-PL when the value group is $\mathbb Q$.

By GJR Definition 7.8, a real-PL function $h$ is LFS at $x$ precisely when its residue class $L_h(x)\in\operatorname{Pic}(\operatorname{red}(X,x))\otimes\mathbb R$ is semiample. Semiample for a real class means a finite nonnegative real combination of semiample integral classes. GJR 7.9 gives its ordinary-local description as a nonnegative real combination of maxima of negative logarithms of analytic units. The prerequisite that $h$ is real-PL is retained. A logarithm of a unit, and any real constant, has zero residue and is LFS; GJR Lemmas 6.17--6.18 prove the residue statement. For constants choose $a\in K^\times$ with $|a|\neq1$ and write the constant as a real multiple of $\log|a|$.

GJR Proposition 7.10 gives a sequence of uniformly convergent LFS functions near an interior point of a semipositive real-PL function. Its conclusion does not say decreasing. Its proof uses $h+\epsilon h'$ on one neighborhood for this one $h$. Applying it separately to an infinite family can change the neighborhood and is not our construction.

GJR 15.2--15.5 provide finite semistable skeletons and integral PL realization for piecewise $(\mathbb Z,\Gamma)$-linear functions. Thu Theorem 2.3.8 supplies the finite separable semistable extension for a strict affinoid smooth curve; over our algebraically closed $K$ that extension is $K$ itself. Thu Lemma 2.2.22 puts finitely many type-2 points in a skeleton, and then makes them vertices, again over $K$. Its part (i) requires type-2 points for the finite separable extension. We never turn an arbitrary type-3 point into a vertex over $K$.

The skeleton slope-degree formula and harmonicity are used in exactly the form checked in GJR Proposition 15.7 and its proof: at an interior vertex the outgoing slope sum is the degree of the residue line bundle, and vanishing of that sum is harmonicity. GRpl Proposition 5.1, and equivalently GRc Theorem 4.3 on a good strict space, identify subharmonic strict real-PL functions on a smooth curve with nef residue classes. The numbering is important: the cached GRpl version calls the curve comparison Proposition 5.1, although GRc cites another version's Proposition 3.1.

Wanner Corollary 5.4 assumes a finite continuous target and either residue field algebraic over a finite field, projective line, or Mumford curve. Its Dini argument is not used for our arbitrary targets. GJR Proposition 15.11 states a decreasing net of weakly smooth functions; it does not state a sequence or an LFS conclusion.

The cached CD PDF numbers its semiample-residue implication as Proposition 17.8.2, not 6.8.2. Its hypotheses there are a separated almost-good space, an integral paralinear metric on an ordinary-locally trivial line bundle, and a positive residue power generated by global sections. We use the directly stated GJR 7.8--7.9 convention, not a mismatched numbered CD proposition or a claimed comparison of smoothness classes.

\section*{3. A strict-domain realization lemma}

Let $Y$ be a compact connected strictly affinoid smooth curve with a semistable model over $K$, skeleton $\Sigma$, and retraction $\tau:Y\to\Sigma$. Suppose a continuous function $F$ on the finite graph $\Sigma$ is affine on every edge of a subdivision whose vertices are type 2. Slopes and function values can be arbitrary real numbers. Then $F\circ\tau$ is strict-domain real-PL on $Y$.

Here is the local proof, which does not round real breakpoints or invoke Thu's all-domain sheafification. At a type-2 vertex $v$, cut each of the finitely many incident edges at a nearby type-2 point. Distances of the cuts from $v$ belong to $\Gamma$. Refine the semistable model at these cuts; the inverse image of this closed star is a compact strictly analytic neighborhood with that star as skeleton. For each incident edge $e$ define $F_e$ to have value zero at $v$, slope one away from $v$ on $e$, and value zero on the other edges. With edge coordinates normalized at their type-2 endpoints, $F_e$ has integral slopes and intercepts in $\Gamma$. GJR 15.4--15.5 realize $F_e\circ\tau$ as an integral PL function on this strict domain. On this star,

\[F\circ\tau=F(v)+\sum_e s_e(F_e\circ\tau),\]

where $s_e\in\mathbb R$ is the outgoing slope. The constant is real-PL by the preceding constant formula. This proves strict real-PL near $v$. At a type-3 point in the interior of an edge, the function is affine in an annulus coordinate, hence is a real multiple of the logarithm of a unit plus a constant on a sufficiently small strict annulus neighborhood. At points off the skeleton it is locally constant in the open discs contracted by $\tau$, so is real-PL there. Real-PL functions form a sheaf in the ordinary topology, GJR 5.2(4), so these local proofs give the assertion on $Y$.

The local star can include a boundary vertex: the PL realization itself needs no positivity at that boundary. For positivity below we always pass to $Y\setminus\partial Y$. This proof permits arbitrary real coefficients but only value-group vertices. It therefore does not contradict GJR Example 5.3.

\section*{4. Finite interpolation on one fixed ordinary open set}

Fix $c\in C$. Choose a compact connected strictly affinoid neighborhood $Y\subset C$ with $c\in U:=Y\setminus\partial Y$. Here $\partial Y$ is the actual analytic boundary. It is a finite type-2 set over algebraically closed $K$. Since $C$ is boundaryless and separated, GRc 2.11(3)--(4) identifies $U$ with the topological interior of $Y$ in $C$, so $U$ is an ordinary open neighborhood. Shrink inside the connected component of $C$ under consideration. The connected affinoid curve $Y$ has nonempty boundary. Indeed a compact boundaryless connected affinoid of positive dimension would, by the maximum principle applied to all $\log|a|$, have the same value of every $a$ at every point, making its affinoid spectrum a singleton, impossible in positive dimension over algebraically closed $K$.

If $g$ is identically $-\infty$ on this component, the sequence of constants $-j$ already solves the local problem. Otherwise $g$ is finite at all type-2 and type-3 points, by Thu Proposition 3.4.10 or GRc Proposition 3.5(9).

For each finite set $S$ of type-2 points in $Y$ containing $\partial Y$, construct $H_S$ as follows. Refine a semistable model of $Y$ so that $S$ lies in the vertex set of its finite skeleton. Assign the value $g(s)$ at each $s\in S$. On each edge use an affine function and at every other vertex impose zero outgoing slope sum. This finite Dirichlet system has a unique solution: its quadratic energy is strictly positive for a nonzero vertex function which is zero on $S$, because every connected graph component meets the boundary in $S$. Pull it back by the retraction. It is continuous and finite on $Y$, equals $g$ on $S$, and is harmonic on $Y\setminus S$ by the slope-degree formula and GJR 15.7. This is the unique harmonic interpolation used in Thu's proof of Theorem 3.4.2. The construction is independent of the model: uniqueness follows from the harmonic maximum principle with the assigned boundary values. Section 3 proves it is strict-domain real-PL.

For completeness the comparison arguments are as follows. The subharmonic maximum principle on each component of $Y\setminus S$ gives $g\leq H_S$ on $Y$. If $S\subset S'$, compare $H_S$ with $H_{S'}$ on $Y\setminus S'$: on its boundary $S'$ the first function is at least $g=H_{S'}$, and both are harmonic there. Hence $H_S\geq H_{S'}$ on $Y$.

The function $H_S|_U$ is subharmonic. Given an affinoid domain $V\subset U$ and a harmonic function $h$ with $h\geq H_S$ on $\partial V$, first $h\geq g$ on $V$ by subharmonicity of $g$. At $S\cap V$ this is $h\geq H_S$. On $V\setminus S$ compare the two harmonic functions using $\partial V\cup(S\cap V)$ as boundary. It follows that $h\geq H_S$ on $V$. This proves the subharmonic test; it is precisely the comparison in Thu's published interpolation proof. Consequently $H_S$ has nef residues on $U$ by the strict real-PL comparison.

Let $w=\inf_S H_S|_U$. The infimum of this decreasing net is subharmonic, or identically $-\infty$, by Thu 3.1.8(ii). It agrees with $g$ at every type-2 point of $U$, because that point can be included in $S$. In particular it is not identically $-\infty$. At a type-3 point take its annulus segment with type-2 points dense on the segment. Both functions are finite convex continuous on the segment, by Thu Proposition 3.4.6 applied to a surrounding semistable domain. They agree on its type-2 points and therefore agree at the type-3 point. Thu Lemma 3.4.1 says that a subharmonic function is determined by its values on $I(U)$, the type-2 and type-3 points; its proof uses shrinking one-Shilov-point domains and upper semicontinuity to recover the other points. Thus $w=g|_U$, including any type-1 poles. This proves convergence for the type-2-only interpolation net without declaring every all-domain lisse function strict real-PL.

\section*{5. One perturbation handles every residue class}

Choose one initial semistable model of $Y$ with skeleton $\Sigma_0$. Let $B$ be the finite set of its interior vertices whose smooth proper residue curves have positive genus. All positive-genus type-2 points of $U$ lie in $B$. To check this assertion, off the initial vertices a point is either on an open annulus or inside a contracted open disc; a type-2 point there has residue curve $\mathbb P^1_{\widetilde K}$. The description follows from the smooth and nodal tubes of a semistable curve. Adding vertices by the blowups in Thu 2.2.22 only adds these genus-zero points. Hence the set $B$ is unchanged when an interpolation skeleton is refined.

On $\Sigma_0$ solve the graph Dirichlet problem with value zero at $\partial Y$ and outgoing slope sum one at each $b\in B$, zero at every other interior vertex. To see existence without any positivity assumption, for unknown interior vertex values $q_v$ solve

\[\sum_{e=vw}\frac{q_w-q_v}{\ell_e}=\mu_v,\qquad \mu_v=\begin{cases}1&v\in B,\\0&v\notin B.\end{cases}\]

Parallel edges are counted separately. The negative of the coefficient matrix has quadratic form $\sum_{e=vw}(a_v-a_w)^2/\ell_e$, with boundary values zero. It is positive definite, as every connected component meets the boundary. Thus the system is invertible. Extend the solution affinely on edges. It is finite and continuous on the compact graph; by the graph maximum principle it is nonpositive. Put

\[Q=q\circ\tau_0+A,\qquad A=\max_{y\in Y}|q(\tau_0(y))|.\]

Then $Q\geq0$ on $Y$. Section 3 proves strict real-PL membership. The outgoing slope sums are nonnegative on $U$, and off the skeleton the function is locally constant. The slope-degree formula and the curve comparison show that $Q$ is subharmonic and has nef residues on $U$. At each $b\in B$ the residue degree is exactly one. It may have negative outgoing slope sums at boundary vertices, which have been removed from $U$; no positivity assertion there is needed.

The required algebraic fact for real residue classes is: on a smooth proper connected curve over an algebraically closed field, every class in $\operatorname{Pic}\otimes\mathbb R$ of strictly positive degree is ample in the real-class sense, hence semiample. Here is the real-coefficient justification. Write the class as $\sum_{i=1}^m a_i L_i$ with integral $L_i$. The integral degrees define a rational linear functional on $\mathbb R^m$. The vector $(a_i)$ lies in its open positive half-space. A small rational simplex in that half-space can be chosen to contain this vector in its interior: approximate the vertices of a surrounding real simplex by rational vectors, preserving containment and positive degree. Express the vector with positive barycentric coefficients. Clearing each rational vertex's denominators expresses the class as a nonnegative real combination of integral positive-degree line bundles. An integral line bundle of positive degree is ample and has a globally generated positive power: for sufficiently large powers its degree is greater than $2g$, and the curve Riemann--Roch/duality criterion applied after subtracting any point gives global generation (degree greater than $2g+1$ separates length-two subschemes as well). This proves the assertion. On $\mathbb P^1$, $\operatorname{Pic}\otimes\mathbb R=\mathbb R[\mathcal O(1)]$, so every nonnegative-degree class is already semiample, including zero.

At an interior type-2 point of a smooth curve, its germ reduction has the smooth proper residue curve as model. This follows also by taking the component at the corresponding semistable vertex in GJR 6.13--6.16. Normalizing and completing that function field gives the unique smooth proper curve; proper birational maps between smooth curves are isomorphisms. Thus representing and testing the residue class on that curve is legitimate. At non-type-2 points the functions under consideration are harmonic locally; their residue class is zero, since the residual function field has transcendence degree zero over the algebraically closed residue field. In particular it is semiample.

For every finite $S$ as above and every $\epsilon>0$, define

\[V_{S,\epsilon}=H_S+\epsilon Q\quad\hbox{on }U.\]

This is a strict real-PL function. Its residue class is the sum of the two residue classes by GJR 6.18. At $b\in B$, it has degree at least $\epsilon$, so the algebraic fact makes it ample. At every other type-2 point its residue curve has genus zero and its degree is nonnegative, so its class is semiample. At non-type-2 points it has zero residue. Therefore every $V_{S,\epsilon}$ is LFS at every point of the same ordinary open set $U$, by GJR 7.8. GJR 7.9 also guarantees actual ordinary-local formulas by logarithmic maxima there. No claim that nef equals semiample on an arbitrary higher-dimensional residue variety is made.

Order the pairs by enlargement of $S$ and decrease of $\epsilon$ to zero. Since $Q\geq0$, this is a decreasing LFS net. Its infimum is $g|_U$, by the preceding interpolation convergence and $\epsilon Q\to0$. This handles arbitrary residue Picard classes, not just torsion ones. For example a non-torsion degree-zero class can prevent a particular harmonic real-PL function from being LFS, but the positive-degree perturbation makes its approximants LFS. Non-torsion $\operatorname{Pic}^0$ is therefore not, by itself, a failure criterion for decreasing local approximation.

\section*{6. Sequence extraction without continuity of the target}

Assume now that $K$ has a countable dense subset. A strictly affinoid algebra over $K$ is a separable Banach space: polynomials in its finitely many Tate generators with coefficients in a countable dense set form a countable dense subset, and a quotient of a separable Banach space is separable. On its spectrum, evaluation of $|a|$ is bounded by the Banach norm and satisfies $\big||a(x)|-|b(x)|\big|\leq\|a-b\|$. Therefore a countable dense collection of algebra elements induces the full topology of its compact Hausdorff spectrum and separates points. The spectrum embeds into a countable product of compact real intervals. It is second countable. In particular $Y$ and its open subset $U$ are second countable; every open subset of $U$ is Lindelof.

For each $t\in\mathbb Q$, let $E_t=\{x\in U:g(x)<t\}$. This is ordinary open by upper semicontinuity, and

\[E_t=\bigcup_S\{x\in U:H_S(x)<t\}.\]

Lindelofness gives a countable subcover for each $t$. List all finite sets selected by these countably many subcovers as $T_1,T_2,\ldots$ and put $S_j=\partial Y\cup T_1\cup\cdots\cup T_j$. If the list is finite, repeat its final set; if it is empty, use $\partial Y$. The $H_{S_j}$ decrease. For any $x$ and any rational $t>g(x)$, the chosen subcover of $E_t$ supplies an interpolant $H_{T_k}(x)<t$, so $H_{S_j}(x)<t$ for all $j\geq k$. Together with $H_{S_j}\geq g$ this proves $H_{S_j}(x)\downarrow g(x)$. If $g(x)=-\infty$, use every rational $t$; the same argument gives convergence to $-\infty$.

Finally set

\[v_j=H_{S_j}+2^{-j}Q\quad\hbox{on the fixed }U.\]

Each $v_j$ is finite LFS by Section 5. Both summands decrease, since $Q\geq0$, so $v_{j+1}\leq v_j$. The limit is $g|_U$. This extraction need not be cofinal in the interpolation net. It is a countable extraction by sublevel covers, not Dini's theorem, and it does not presume a continuous target. Local uniform convergence is claimed only where the target is finite continuous on a compact set, in which case the usual compact decreasing-continuous-function argument applies.

\section*{7. Actual projective bundles and proper factors}

First consider a nonempty geometrically connected proper algebraic variety $P$ over $K$. Over every analytic extension field $F/K$, the analytification of $P_F$ is compact, boundaryless and connected. Properness is preserved by scalar extension and analytification, Berkovich Proposition 2.6.9; GJR 2.3 records goodness, strictness and empty boundary for algebraic analytifications. For connectedness one can check an actual published theorem rather than assume it: Berkovich Theorem 7.1.1 identifies compactly supported cohomology of a compactifiable algebraic scheme with that of its analytification for an abelian torsion sheaf. Apply it in degree zero to a constant nonzero finite sheaf on $P_F$. Since $P_F$ is proper, compactly supported degree zero equals degree zero. Its algebraic connectedness gives only constant sections. A disconnected analytic space would have a nonconstant locally constant section, contradicting that comparison. This works for every complete $F$ and does not assume $F$ algebraically closed. In particular it applies to projective spaces. Because $K$ is algebraically closed, a nonempty $P$ has a $K$-rational closed point, giving a product section.

Let $u$ be any classical psh function on a saturated open subset of one of the stated total spaces. For $c\in C$, its pullback to the fibre over $\mathcal H(c)$ is classical psh by GRc 4.2 scalar extension and functoriality. That fibre is a compact connected boundaryless analytic space: it is $\mathbb P^r_{\mathcal H(c)}{}^{\mathrm{an}}$ for a projective bundle and $P_{\mathcal H(c)}^{\mathrm{an}}$ for a product. Upper semicontinuity gives a maximum unless all values are $-\infty$, in which case constancy is immediate. Otherwise the maximum is finite. By GRc Theorem 6.1 its maximum locus is open, and by upper semicontinuity it is closed. Connectedness makes that locus the whole fibre. Hence $u$ is constant on each fibre.

Define $g(c)$ to be this fibre value. On a projective-bundle trivializing ordinary open subset, the constant projective coordinate $[1:0:\cdots:0]$ is an actual analytic section $s$, and $g=u\circ s$. On a product use the fixed rational point of $P$. Thus $g$ is upper semicontinuous and classically psh on each such base open set by pullback; ordinary sheafiness, GRc 4.2(1), gives it on the base. The section values on overlaps agree because they lie in the same fibre. This proves $u=g\circ f$ and is a field-of-definition descent argument, not an assumption that a function on the total space already came from the base.

Apply Sections 4--6 to $g$, choosing $Y$ inside one base neighborhood when necessary. On $f^{-1}(U)$ put $w_j=v_j\circ f$ in the countable case, or pull back the net otherwise. This set is one fixed ordinary open neighborhood of the entire fibre. Pullback preserves strict real-PL by GJR 5.4(1). It preserves LFS directly: pull back each local formula of GJR 7.9; every analytic unit remains a unit, and the same finite maximum formula gives the pulled function. The functions $w_j$ are finite, decreasing, and converge pointwise to $u$.

Projective transition matrices do not alter any $w_j$, since it is the pullback of one globally defined function on $U$. On overlaps of fibre charts the functions agree exactly, not just in their limit. Different choices of base neighborhoods may give different regularizing sequences, but their restrictions have the same target on the overlap; the local assertion does not demand their global gluing.

For every complete analytic field extension $L/K$, the same functions pull back to $f_L^{-1}(U_L)$, where $U_L$ is the inverse image ordinary open set. By GJR 5.4(2) they are strict real-PL, and the logarithmic formulas prove LFS there. Evaluating at a point above $x$ shows the inequalities and the limiting value are unchanged: the limit is $u\circ\pi_{L/K}$. All construction data and functions were already over $K$; no norm or averaging descent from an unspecified extension is needed. This proves compatibility for pullbacks of the original target. It does not assert that every new psh function over an arbitrary nonseparable $L$ has a sequence: for that fresh target the countably-dense field hypothesis must be checked again.

\section*{8. A sharp boundary sequence obstruction}

Let $D=\mathcal M(K\langle T\rangle)$ and let $\zeta$ be its Gauss point. Its actual analytic boundary is $\{\zeta\}$. Define $b(\zeta)=1$ and $b(x)=0$ for $x\neq\zeta$. It is upper semicontinuous, since its positive upper level sets are either $D$, $\{\zeta\}$ or empty. On $\operatorname{Int}(D)=D\setminus\{\zeta\}$ it is the constant zero and is subharmonic. GRc Proposition 3.5(3) extends subharmonicity from the interior to $D$ for this upper semicontinuous function; GRc 4.1.2 makes it classical psh. Its pullback after every analytic extension is therefore in the same classical class. This argument uses a genuine analytic boundary; the corresponding interior spike on a boundaryless disc is not asserted to be psh.

For algebraically closed $K$, the components of $D\setminus\{\zeta\}$ are the open residue discs $D_a$, $a\in\widetilde K$. Every ordinary neighborhood $N$ of $\zeta$ contains all of $D_a$ for all but finitely many $a$. Indeed $D\setminus N$ is compact and is covered by the disjoint open sets $D_a$, so a finite subcover suffices.

Suppose continuous functions $h_j$ on one ordinary neighborhood $N_0$ of $\zeta$ decrease pointwise to $b$. For each $j,m\geq1$, continuity supplies a neighborhood $N_{j,m}\subset N_0$ of $\zeta$ on which $|h_j(x)-h_j(\zeta)|<1/m$. Each of these neighborhoods excludes only finitely many complete residue discs. If $\widetilde K$ is uncountable, choose $a$ outside the countable union of those finite exceptional sets. Then $D_a\subset N_{j,m}$ for every $j,m$, so every $x\in D_a$ satisfies $h_j(x)=h_j(\zeta)$ for all $j$. Its limiting value must be one, contradicting $b(x)=0$. No continuous decreasing sequence, hence no LFS sequence, exists on $N_0$. In fact this continuity argument forbids any pointwise convergent sequence with this prescribed limit. The dimension-$r+1$ example is $D\times\mathbb P^r_K{}^{\mathrm{an}}$ with target pulled back from $b$. A fixed rational projective point gives a section. Restricting a putative local LFS sequence near $(\zeta,p)$ to this section gives the forbidden continuous sequence on a neighborhood of $\zeta$.

For sharpness for this particular function, suppose instead that $\widetilde K$ is countable, enumerate its elements without repetitions as $a_1,a_2,\ldots$, and choose lifts $t_i\in K^\circ$. Pick $c_0\in K^\times$ with $0<|c_0|<1$, set $\alpha=-\log|c_0|>0$, and for each $j$ choose $d_j\in K$ with $d_j^j=c_0$. Define, with the maximum finite even at a zero of $T-t_i$,

\[b_j=1+\sum_{i=1}^j\frac{j}{\alpha}\max\{\log|T-t_i|,\log|d_j|\}\quad\hbox{on }D.\]

Each summand is strict real-PL. Its bend is at the value-group radius $|d_j|$, with domains $|T-t_i|\leq|d_j|$ and $|T-t_i|\geq|d_j|$. On the first domain it is constant; on the second $T-t_i$ is a unit. Thus this is a strict affinoid G-cover, unlike an arbitrary real-radius subdivision. For LFS membership near a zero the term is locally constant. Away from its zero the unscaled maximum is $\max\{-\log|(T-t_i)^{-1}|,-\log|d_j^{-1}|\}$, the GJR 7.9 unit formula. The coefficient is nonnegative; sums and constants have semiample residues, so $b_j$ is LFS everywhere on $D$.

At $\zeta$ every summand is zero, so $b_j(\zeta)=1$. At $x\in D_{a_k}$ all terms except the $k$th are zero, and the $k$th contribution, when included, is $\max\{(j/\alpha)\log|T-t_k|(x),-1\}$. Therefore $0\leq b_j\leq1$, $b_{j+1}\leq b_j$, and at each such $x$ it eventually equals zero. The same is true at $T=t_k$ by the constant branch formula. Hence $b_j\downarrow b$ on the fixed disc $D$. All roots, lifts and strict domains were defined over $K$; pulling back the formula proves LFS, inequalities and convergence to the pulled target over every analytic extension.

Even with uncountable residue field there is a decreasing LFS net for $b$: use finite subsets of residue classes in the displayed sum and a positive integer truncation parameter $n$, with $d_n^n=c_0$, and order by enlargement of the subset and increase of $n$. At each residue disc its class eventually enters the subset and its term tends to $-1$. Thus the obstruction is exactly to sequences for this function, not to nets or to LFS approximation in general.

\section*{9. Remaining original obligations and evidence limits}

The original P02.2 is not solved. The argument gives a sufficient saturated proper-fibre sector, not necessary and sufficient geometry for all higher-dimensional spaces or for every local psh germ. General arbitrary psh functions on nonsaturated higher-dimensional neighborhoods need not descend to curves. No construction is supplied for them. Boundaryless higher-dimensional residue varieties can have nef classes outside the semiample cone; the curve positive-degree argument does not extend to these varieties. We have not proved that this algebraic distinction obstructs approximation, because a perturbation can sometimes remove it. The boundary example gives a sharp residue-cardinality criterion for its specific sequence problem; it does not characterize the full failure range and does not obstruct nets. We have not removed algebraic closure from the positive construction or proved a descent of its newly chosen interpolants and perturbation over general nonclosed fields. Arbitrary targets over nonmetrizable fields receive a net here; a general sequence theorem is not asserted.

Rerun the lightweight evidence with {\small\texttt{python collect\_sources.py}}, {\small\texttt{python finite\_diagnostics.py}}, {\small\texttt{python build\_documents.py}}, and {\small\texttt{python verify\_artifacts.py}} in this workspace. The standalone {\small\texttt{P02.2-paper.tex}} embeds the entire argument, and the build program runs the existing TeX installation twice. {\small\texttt{source-extracts/source-provenance.json}} records exact cached hashes, retained literal text and actual PDF pages, numbering differences, and application checks. {\small\texttt{finite-diagnostics.json}} checks only finite rational graph examples and a finite interpolation grid. The original input is retained byte-identically as {\small\texttt{input-original.json}}; the parsed assignment and task are in {\small\texttt{assignment-original.txt}} and {\small\texttt{task-original.json}}. {\small\texttt{status.json}} uses only the five requested keys.

Source hashes establish the identity of the inspected files. Finite diagnostics, successful TeX compilation and artifact validation do not certify any infinite or mathematical claim. The mathematical proofs are the arguments above and the exactly identified published foundations; the whole submission remains an unverified draft with candidate false.
\end{document}
